\documentclass[11pt]{article}

% ---- packages ----
\usepackage[margin=1in]{geometry}
\usepackage{amsmath,amssymb}
\usepackage{booktabs}
\usepackage{array}
\usepackage{xcolor}
\usepackage[hidelinks]{hyperref}
\usepackage{enumitem}
\usepackage{microtype}
\usepackage[htt]{hyphenat}
\usepackage{seqsplit}
\usepackage[ruled,vlined,linesnumbered]{algorithm2e}

% ---- styling ----
\definecolor{unityblue}{HTML}{2C5F9E}
\definecolor{pyorange}{HTML}{C9560B}
\definecolor{flagred}{HTML}{B3261E}
\newcommand{\code}[1]{{\texttt{\small #1}}}
\newcommand{\unity}[1]{{\color{unityblue}\texttt{\small #1}}}
\newcommand{\py}[1]{{\color{pyorange}\texttt{\small #1}}}
\setlength{\emergencystretch}{3em}  % avoid overfull boxes from long unbreakable mono tokens

\title{\textbf{The Greedy and Rules-Based Reference Policies\\ in the ARC Disaster-Response Environment}}
\author{ARC\_Game RL --- technical reference}
\date{\today}

\begin{document}
\maketitle

\begin{abstract}
\noindent This note documents the two non-learning, API-free reference policies used as baselines:
the \emph{greedy} policy (\py{greedy\_decision}) and the \emph{rules-based} policy
(\py{potential\_decision}), both in \py{benchmark\_models.py}. They are presented separately:
greedy is a purely reactive policy that never invests capital; rules-based reuses greedy's reactive
core and adds a thin capital-investment layer. A final section records their \textbf{known
limitations} --- both are deliberately simple heuristics that under-use the action space and do no
forward planning, which is precisely the headroom a learned policy is meant to exploit. Reward
terms referenced below are defined in the companion document, \emph{Reward Scoring and Data
Provenance}.
\end{abstract}

\tableofcontents
\vspace{1em}

\paragraph{How a turn is executed.} Both policies return, each round, (i)~a set of task-choice
selections $\{(\text{taskId}, \text{choiceId})\}$ and (ii)~a list of game-action indices. The
environment executes \emph{all} of them: \py{env.step()} takes a comma-separated list of action
indices, so a turn may contain arbitrarily many actions of any mix (multiple builds, multiple
hires, builds and hires together). Any per-turn limits below are \emph{self-imposed by the policy},
not imposed by the environment.

% =====================================================================
\section{The Greedy Policy}
% =====================================================================

The greedy policy is purely \textbf{reactive and myopic}. Each round it does exactly two things:
\begin{enumerate}[leftmargin=1.5em, itemsep=1pt]
\item answers open tasks by selecting, for \emph{every} task, the single choice that maximizes an
immediate value function (if any choice scores positively); and
\item assigns idle workers (at zero cost) to buildings that are short-staffed.
\end{enumerate}
It \textbf{never builds and never hires} --- it performs no capital investment whatsoever. It
simply spends the starting infrastructure and workforce as efficiently as the current tasks allow.

\subsection{Task-choice value function}
For an active task $t$, each available choice $c$ exposes a budget impact $b(c)$ (positive =
funding gained, negative = a cost) and a satisfaction impact $s(c)$. With the demand indicator
\[
D(t) = \begin{cases}1 & \text{taskType} \in \{\text{Demand},\,\text{Emergency}\}\\ 0 & \text{otherwise,}\end{cases}
\]
the value of a choice is
\begin{equation}
v(c) =
\begin{cases}
\dfrac{b(c)}{10^4} + 10^{-2}\,s(c), & b(c) > 0 \quad\text{(funding choice)}\\[10pt]
\mathbf{1}\!\big[\alpha(c)\wedge D(t)\big] + 10^{-2}\,s(c) - w_{c}\,\mathrm{cost}(c), & b(c) \le 0 \quad\text{(acting / waiting)}
\end{cases}
\label{eq:vfun}
\end{equation}
where $\mathrm{cost}(c) = -b(c) \ge 0$, $w_c = 0.0002$ (matching the reward's cost scale), and the
\emph{acting} predicate is $\alpha(c) = [\mathrm{cost}(c) > 0] \vee [s(c) \ge 10]$. For each task
the policy selects $c^\star = \arg\max_c v(c)$, but only records it if $v(c^\star) > 0$. The
$+1$ term rewards actually fulfilling a genuine demand/emergency; funding offers are ranked by
dollars; per-unit cost is subtracted. This favours the \emph{reliable, immediate} fulfilment
options over ``free'' deferred ones that often fail to complete.

\paragraph{Many tasks, many free answers --- in one turn.} The value function is evaluated
\emph{independently for every open task}, and a selection is emitted for each task whose best choice
clears $v>0$. There is no cap on how many tasks are answered per turn: if several relocation or food
tasks are open --- including several that can be answered at no cost (a free choice carrying any
positive satisfaction signal scores $v = 10^{-2}s > 0$) --- the policy answers \emph{all of them}
in the same round.

\subsection{Worker assignment}
The greedy policy also issues zero-cost worker assignments: it groups the enumerated
\unity{worker\_assignment} actions by target building and, preferring trained workers and larger
batches, fills \emph{each} short-staffed building once from the free worker pool.

\begin{algorithm}[h]
\DontPrintSemicolon
\SetKwInOut{Input}{input}\SetKwInOut{Output}{output}
\Input{game state; enumerated valid actions $\mathcal{A}$}
\Output{choices $C$; action indices $X$}
$C \leftarrow \emptyset$;\quad $X \leftarrow \emptyset$\;
\ForEach{active task $t$ with choices}{
  $c^\star \leftarrow \arg\max_{c \in t}\, v(c)$ \tcp*{value function, Eq.~\eqref{eq:vfun}}
  \If{$v(c^\star) > 0$}{ add $(t, c^\star)$ to $C$ \tcp*{every answerable task, no per-turn cap}}
}
$f_{\mathrm{tr}}, f_{\mathrm{un}} \leftarrow$ free trained, free untrained workers\;
group \unity{worker\_assignment} actions by target building\;
\ForEach{building $b$ that needs workers}{
  \ForEach{candidate $a$ of $b$, ordered by (trained first, larger quantity)}{
    \If{$0 < \mathrm{qty}(a) \le$ available workers of $a$'s type}{
      add index of $a$ to $X$; decrement that worker pool; \textbf{break}\;
    }
  }
}
\Return $(C, X)$ \tcp*{note: no construction, no hiring --- ever}
\caption{\textsc{GreedyPolicy} (\py{greedy\_decision})}
\label{alg:greedy}
\end{algorithm}

% =====================================================================
\section{The Rules-Based Policy}
% =====================================================================

The rules-based policy reuses the greedy policy verbatim as its \textbf{reactive core}, then
\emph{appends} a small amount of capital investment --- the one thing greedy never does. The name
mirrors potential-based shaping: it spends now against a hand-crafted \emph{state potential}
(durable shelter / kitchen / casework capacity and staffed workers) expected to pay off later
through the satisfaction and worker-use reward terms.

\subsection{State statistics and the demand anchor}
From the current facility list it computes
\begin{align}
P &= \textstyle\sum_{\text{Community }f} \text{currentPopulation}(f) \quad\text{(demand anchor; default 120 if unknown)}, \\
\text{shelterCap} &= \textstyle\sum_{\text{Shelter }f} \text{populationCapacity}(f), \quad
n_{\text{kit}} = \#\{\text{Kitchen}\}, \quad
n_{\text{cw}} = \#\{\text{CaseworkSite}\}.
\end{align}
$P$ is the \emph{current} total displaced population --- a \textbf{static anchor, not a forecast}.
The policy sizes shelter capacity toward $P$ so people can eventually be housed for free in
shelters rather than in the recurring-cost motel; it does not anticipate the timing or size of
future demand waves (see~\S\ref{sec:limits}).

\subsection{Capital investment: build then hire}
After taking greedy's $(C,X)$, the overlay appends \textbf{at most one construction} and
\textbf{at most one hire}, both gated on horizon and a budget floor:

\paragraph{Build (one building per turn).} If $\text{roundsLeft} \ge H_{\min}$ and
$\text{budget} \ge B_{\mathrm{res}}$, choose a single target by strict priority and commit it only
if it leaves the floor intact:
\begin{equation}
\text{target} =
\begin{cases}
\textsc{CaseworkSite}, & n_{\text{cw}} < 1 \ \wedge\ \text{round} \ge R_{\text{cw}}\\
\textsc{Shelter}, & \text{else if } \text{shelterCap} < P\\
\textsc{Kitchen}, & \text{else if } n_{\text{kit}} < K_{\text{tgt}}\\
\text{none}, & \text{otherwise,}
\end{cases}
\end{equation}
and the chosen target is committed only if $\mathrm{cost}(\text{target}) \le \text{budget} - B_{\mathrm{res}}$.
The casework site is built first ($R_{\text{cw}}=0$): a site is staffed only if free workers exist
when it is built, so claiming its workers early is the only way it stays operational. \emph{Site
selection is geographically blind} --- among the available sites for the chosen type it takes the
cheapest, but since every construction shares one fixed cost this reduces to the
\emph{first-enumerated} available site (see~\S\ref{sec:limits}).

\paragraph{Hire (one worker per turn).} Let the understaffing be
\begin{equation}
\text{need} = \sum_{f:\,\text{NeedWorker}} \max\big(0,\ \text{requiredWorkforce}(f) - \text{assignedWorkforce}(f)\big).
\end{equation}
A single \unity{hire\_untrained} action is appended when $\text{need} > f_{\mathrm{tr}} +
f_{\mathrm{un}}$ and $\text{budget} \ge B_{\mathrm{res}}$.

\begin{algorithm}[h]
\DontPrintSemicolon
\SetKwInOut{Input}{input}\SetKwInOut{Output}{output}
\Input{game state; round $r$; total rounds $T$}
\Output{choices $C$; action indices $X$}
$(C, X) \leftarrow \textsc{GreedyPolicy}(\text{state})$ \tcp*{reactive core, Algorithm~\ref{alg:greedy}}
compute $P,\ \text{shelterCap},\ n_{\text{kit}},\ n_{\text{cw}}$;\quad
$\text{budget}$, $\text{roundsLeft} \leftarrow \max(0,T-r)$\;
\BlankLine
\If{$\mathrm{roundsLeft} \ge H_{\min}$ \textbf{and} $\mathrm{budget} \ge B_{\mathrm{res}}$}{
  $\text{target} \leftarrow$ first applicable: \textsc{CaseworkSite} $(n_{\text{cw}}{<}1 \wedge r{\ge}R_{\text{cw}})$,
  \textsc{Shelter} $(\text{shelterCap}{<}P)$, \textsc{Kitchen} $(n_{\text{kit}}{<}K_{\text{tgt}})$\;
  \If{target exists \textbf{and} $\mathrm{cost}(\text{target}) \le \mathrm{budget} - B_{\mathrm{res}}$}{
    append \emph{one} build action (first available site) to $X$;\quad $\text{budget} \mathrel{-}= \mathrm{cost}$\;
  }
}
\If{$\mathrm{need} > f_{\mathrm{tr}} + f_{\mathrm{un}}$ \textbf{and} $\mathrm{budget} \ge B_{\mathrm{res}}$}{
  append \emph{one} \unity{hire\_untrained} action to $X$\;
}
\Return $(C, X)$\;
\caption{\textsc{RulesBasedPolicy} (\py{potential\_decision}, baseline mode)}
\label{alg:rules}
\end{algorithm}

\begin{table}[h]
\centering
\footnotesize
\renewcommand{\arraystretch}{1.2}
\begin{tabular}{>{\raggedright\arraybackslash}p{3.6cm} l >{\raggedright\arraybackslash}p{7.8cm}}
\toprule
\textbf{Constant} & \textbf{Default} & \textbf{Role} \\
\midrule
$H_{\min}$ (\py{\_POT\_MIN\_HORIZON}) & $4$ & Don't build with fewer rounds left --- can't amortize. \\
$B_{\mathrm{res}}$ (\py{\_POT\_BUDGET\_RESERVE}) & $1500$ & Floor below which no discretionary build/hire occurs. \\
$R_{\text{cw}}$ (\py{\_POT\_CASEWORK\_BUILD\_ROUND}) & $0$ & Earliest round to build the casework site (build-first). \\
$K_{\text{tgt}}$ (\py{\_POT\_KITCHEN\_TARGET}) & $2$ & Operational kitchens to aim for (food + employment). \\
$\theta$ (\py{\_POT\_SHELTER\_COVERAGE}) & $\infty$ & Free-shelter re-routing threshold; $\infty$ disables it. \\
\py{\_POT\_MODE} & \code{baseline} & Selects this policy vs.\ the \code{demandsupply} variant. \\
\bottomrule
\end{tabular}
\caption{Rules-based policy constants (\py{benchmark\_models.py}).}
\label{tab:rbconst}
\end{table}

% =====================================================================
\section{Known Limitations (Flaws of these Baselines)}
\label{sec:limits}
% =====================================================================

Both policies are deliberately simple. The environment permits far richer behavior than either uses;
the gaps below are recorded as \textcolor{flagred}{\textbf{flaws}}, not features, and are the headroom
a learned policy should capture.

\begin{enumerate}[leftmargin=1.6em, itemsep=4pt]
\item \textcolor{flagred}{\textbf{Severe under-use of the multi-action turn.}} The environment
executes a \emph{list} of actions per turn, so a policy may build several buildings, hire several
workers, and do both at once. Greedy performs \emph{zero} capital actions; rules-based caps itself
at \emph{one} build and \emph{one} hire per turn. Capital therefore ramps at roughly one building
per round, which --- combined with the $\sim$1-day construction lag --- badly bottlenecks early
capacity even when budget and multiple free sites are available. A sound policy should build and
hire \emph{multiple} units in a single turn when warranted.

\item \textcolor{flagred}{\textbf{No demand forecast / no forward capital investment.}} Demand
arrives in waves over the episode, but neither policy models this. Greedy is purely reactive; the
rules-based anchor $P$ is the \emph{current} population, used only to size shelter capacity. Neither
invests \emph{ahead} of an expected future wave to capture later reward. A sound policy should
forecast upcoming demand and pre-build/pre-staff capacity against it.

\item \textcolor{flagred}{\textbf{Geographically blind site selection.}} Each candidate site
carries a world \code{position}, communities have positions, and flooding can block relocation
routes --- yet site selection ignores all of it. Because every construction shares one fixed cost,
``cheapest site'' degenerates to the first-enumerated available site. The policy never prefers
shelters near the communities they serve, never avoids flood-exposed or route-blockable sites, and
the \code{position} field is in fact dropped by the action enumerator before the policy ever sees it.

\item \textcolor{flagred}{\textbf{The budget reserve is dead capital.}} $B_{\mathrm{res}}=1500$ is a
hard floor that discretionary build/hire can never cross, so it is hoarded and never deployed for
the capacity it ostensibly protects. Worse, reactive task-choice spending is \emph{not} gated by it
and can still drive the budget below the floor --- so the reserve neither funds infrastructure when
needed nor guarantees solvency. A sound policy should be willing to deploy reserves when the
expected return justifies it.
\end{enumerate}

\paragraph{Alternate mode.} Setting \py{\_POT\_MODE}\code{=demandsupply} swaps the overlay for a
variant that routes each relocation into staffed shelter space via the reliable immediate
evacuation while beds last, spilling to the motel only when shelters are full, and sizes shelter
capacity to $P$ via a coverage fraction \py{\_POT\_DS\_COVERAGE}. It shares all four limitations
above; the baseline mode is the default reference policy.

\end{document}
